\documentclass{subfiles}
\begin{document}\label{sec:diffie-hellman}
    Born by the collaboration between \citeauthor{DiffieH1976}, 
        the Diffie-Hellman key exchange protocol was the first public key exchange algorithm.
        Roughuly speaking the algorithm is the following: 
        we select a prime number \(p\) and consider and choose an element \(\gamma \in \gf_{p}\),
        both parties, in our case \(\A \text{ and } \B\), 
        choose a private key, let these be \(a\) for \(\A\) and \(b\) for \(\B\),
        and compute the public key by taking the \(a\)-th (\(b\)-th resp.) power of \(\gamma\).

    A more detailed description of the algorithm is shown in \autoref{proc:DH-exchange}.
        \subfile{../TikZ/Procedure - Diffie-Hellman exchange}

    As stated, Diffie-Hellman is more a public key exchange protocol rather then an encryption schema;
        though, as described in the next section, when combined with Vernam cyphering algorithm,
        it allows to build a solid public key encryption schema known as the \emph{ElGamal} encryption
        schema.
\end{document}
